\slajdid{02-s044}
\begin{frame}[label=02-s044]
Da bi dobili minimalnu funkciju u obliku proizvoda zbirova prekrivamo sve logičke nule površinama najvišeg reda, koristeći se i sa stanjem b

\begin{columns}[c,onlytextwidth]
\begin{column}{.37\textwidth}\centering
% Original prikazuje samo grupu celog reda DC=11; ostale grupe se ne dopisuju.
\begingroup
\fontsize{12}{13.3}\selectfont
\renewcommand{\small}{\fontsize{12}{13.3}\selectfont}
\scalebox{.75}{%
\begin{karnaugh-map}(label=middle,implicantcolors={DEplava})[4][4][1][$A$][$B$][$C$][$D$]
\minterms{0,4,5,8,10}
\terms{1}{b}
\maxterms{2,3,6,7,9,11,12,13,14,15}
\implicant{12}{14}
\end{karnaugh-map}}
\endgroup

\end{column}
\begin{column}{.60\textwidth}
Pošto suštinski površine predstavljaju objedinjavanje proizvoda kod kojih je moguće uraditi sažimanje, u zbirovima ostaju komplementi literala koji se \alert{\underline{ne menjaju}} na povšini.
{\fontsize{10}{12}\selectfont\[
F=(D+\bar B)(\bar D+\bar C)(\bar D+\bar A)
\]}
\end{column}
\end{columns}
\end{frame}
